% Folder 79, batch 14 — archivists' pages 261–280 (author's pp. 134–138 on 262–270, then loose sheets).
% Pass: Opus 5.5 (claude-opus-5-5), 2026-10-03 — first pass, unchecked against the pages by a human.
%
% Notation fixed for this batch: the symmetric group of S is \mathfrak{S}_S;
% the group he writes Z with a tilde is \widetilde{Z}; the automorphism of H
% induced by g is \bar{g}; star sets are \mathrm{Et}. Pages 271 and 274 are
% blank and are absent.
\documentclass[11pt,a4paper]{article}
\input{../preamble/grothendieck}

\folder{79}
\batch{14}
\pages{261}{280}
\foldertitle{2-polyèdres réguliers triangulés, et modules libres de rangs 2 sur les anneaux $\Lambda=\mathbb{Z}/n\mathbb{Z}$ et $\Lambda=\mathbb{Z}$ : notes manuscrites (s.d.).}
\dating{[à partir de 1980]}
\watermark{Édition de démonstration}

\begin{document}

\page{261}
\note{la page continue sans rupture un argument commencé avant ce lot : les renvois (119), (135) et les données $H$, $U_0$, $T_H$, $\sigma$, $Z$, $\widetilde{Z}$, $T_0$, $T_1$, $T_G$ viennent des pages précédentes.}
\marginal{$U_0$, $T_H$ ss-groupes de $H$}
i.e.\ $(H$, $U_0\subset H$, \struck{\ill{}}, $T_H\subset H$, $\sigma\in H)$ satisfaisant
les conditions
\[
(135)\qquad
\begin{cases}
\sigma^{2}=1\\
U_0\cap U_1=\{1\} \quad \text{où } U_1=\sigma(U_0)\\
\sigma(t)=t^{-1} \quad \forall t\in T_H\\
T_H \text{ normalise } U_0\\
U_0\times T_H\times U_1\to H,\ (u_0,t,u_1)\mapsto u_0tu_1 \ \text{injectif}
\end{cases}
\]
(où $U_1=\sigma(U_0)$) ---

il reste : expliciter, en termes de ces données, le groupe $\widetilde{Z}$, les
deux homomorphismes
\[
\widetilde{Z}\ \underset{c'}{\overset{c}{\rightleftarrows}}\ Z
\]
\struck{\ill{}} (où $Z$ est une copie conforme de $T_H$ via l'isom.\
$u\colon Z\xrightarrow{\sim}T_H$) et enfin l'opération de $\widetilde{Z}$ sur $H$, via
\[
\tilde\lambda\longmapsto \operatorname{int}(u_0(\tilde\lambda))|H
\ \overset{\text{déf}}{=}\ \Theta_0(\tilde\lambda)
\]
\marginal{donc à l'opération de $T_0$ sur $H$}
\struck{\ill{}} ou, au choix, via
\[
\tilde\lambda\longmapsto \operatorname{int}(u_1(\tilde\lambda))|H
\ \overset{\text{déf}}{=}\ \Theta_1(\tilde\lambda)
\]
\marginal{i.e.\ l'opération de $T_1$ sur $H$.}
ces deux opérations étant \add{en effet} reliées par la formule --- déduite de
\[
(119)\qquad u_1(\tilde\lambda)=u_0(\tilde\lambda)\,u(c(\tilde\lambda))
\]
qui donne
\[
(136)\qquad \Theta_1(\tilde\lambda)=\Theta_0(\tilde\lambda)\,\operatorname{int}(u(c(\tilde\lambda))) ,
\]
--- \struck{telle} de sorte que $\Theta_0(\tilde\lambda)$, $\Theta_1(\tilde\lambda)$ se déterminent
mutuellement. Il revient aussi au \uncertain{même} de \struck{de} demander, \add{plus
synthétiquement,} la construction du groupe $T_G$, comme groupe opérant sur $H$,
\uncertain{et dans} $T_G$ les \struck{deux} \add{trois} sous-groupes $T_H\hookrightarrow T_G$ (déjà connu :
déterminé), $T_0$ et $T_1$, \struck{$\widetilde{Z}$ étant} \add{\uncertain{que}} \uncertain{alors} défini \emph{de façon}
\struck{\ill{}} symétrique comme étant $T_G/T_H$, et l'homomorphisme $Z\xrightarrow{c'}\widetilde{Z}$
étant décrit comme le composé de

\page{262}
\note{p.~134 de l'auteur.}
\[
Z\xrightarrow[\sim]{u}T_H\hookrightarrow T_G\xrightarrow{\text{can}}\widetilde{Z}.
\]
Il faudra de plus, encore, décrire l'hom.\ $c\colon\widetilde{Z}\to Z$.

Rappelons-nous que l'on a
\[
(137)\qquad
\begin{cases}
T_0=\left\{g\in\mathfrak{S}_S \;\middle|\;
\begin{array}{l}
g \text{ normalise } H\subset\mathfrak{S}_S\\
g \text{ centralise } Z\subset\mathfrak{S}_S\\
gs_0=s_0\\
gt_0\in Z.t_0
\end{array}\right\}\\[4ex]
T_1=\left\{g\in\mathfrak{S}_S \;\middle|\;
\begin{array}{l}
g \text{ normalise } H\\
g \text{ centralise } Z\\
gt_0=t_0\\
gs_0\in Z.s_0
\end{array}\right\}
\end{cases}
\]
\note{dans $T_1$, le mot « centralise » est récrit sur lui-même et se lit mal. Après l'accolade de $T_0$ un signe $\Leftrightarrow$, sans suite ; à droite, d'une écriture plus petite :}
\[
T_G=\left\{g\in\mathfrak{S}_S \;\middle|\;
\begin{array}{l}
g \text{ norma. } H\\
g \text{ centr. } Z\\
gs_0\in Zs_0\\
gt_0\in Zt_0
\end{array}\right\}.
\]

Or pour un espace homogène $S$ sous un groupe $H$ \add{$\subset\mathfrak{S}_S$}, pointé par une
« origine $s_0$ », les groupe des $g\in\mathfrak{S}_S$ qui normalisent $H$ et fixent $s_0$,
i.e.\ des automorphismes de la situation
\[
(S,\ H\subset\mathfrak{S}_S,\ s_0\in S)\longleftrightarrow(H,\ U_0=H_{s_0}),
\]
est can.\ isomorphe au groupe des automorphismes \add{$\bar g$} de $H$ qui
\struck{stabilisent} normalisent $U_0$. \struck{\ill{}} Pour un \struck{\ill{}} $g$, dire qu'il
\struck{centralise $Z\subset\mathfrak{S}_S$, défini} \add{normalise un sous-groupe}
$Z\subset\mathfrak{S}_S$ commutant à $H$, \add{défini} par un sous-groupe
$\underline{B_{0H}}\supset U_0$ du normalisateur de $U_0$, \ill{} signifie que
l'automorphisme correspondant $\bar g$ de $H$ normalise $B_{0H}$. --- et
l'automorphisme $\operatorname{int}(\bar g)|Z$ \struck{\ill{}} s'identifie alors à :
$\operatorname{int}(\bar g)|(B_{0H}/U_0)$.

Donc dire que $g$ \emph{centralise} $Z$, signifie que $\bar g$ normalise $B_{0H}$ et
induit l'identité sur $B_{0H}/U_0$. Il reste : exprimer la condition $gt_0\in Z.t_0$
en termes de $\bar g$, or \struck{$g$} $t_0$ correspond à $\sigma.U_0$, $gt_0$ à
$\bar g(\sigma).U_0$, et \struck{\ill{} qu'il existe} \add{$\lambda t_0$ $(\lambda\in Z)$ à
\uncertain{correspond}}

\page{263}
\struck{$\lambda\in T_H$ tel que $\tilde\lambda\,u(\sigma)\,U_0=$}
\uncertain{prend} : $\sigma.u(\lambda)U_0$, donc la condition
\[
gt_0=\lambda t_0
\]
s'écrit aussi
\[
(138)\qquad \bar g(\sigma)U_0=\sigma.u(\lambda)U_0 \quad \text{pour } \lambda\in Z \text{ convenable}
\]
(\uncertain{automatiquement} unique).

En fait, cette condition implique déjà la condition \ill{}, plus forte, que $\bar g$
\add{(normalise $U_0$ et)} \emph{centralise} $T_H$ \uncertain{lui-même}, et que de plus on ait
\[
(139)\qquad \bar g(\sigma)=\sigma\,u(\lambda) \quad \text{pour } \lambda\in Z \text{ convenable,}
\]
soit encore
\[
(140)\qquad \bar g(\sigma)\in\sigma.T_H ,
\]
\struck{puisque} \add{En effet} en termes de $g\in G\subset\mathfrak{S}_S$ \add{($T_0\subset G$)},
\struck{cette relation} \add{on sait} \uncertain{bien} que $g$ centralise $T_H$, et (137) s'écrit
\struck{s'écrit}
\[
g\sigma g^{-1}=\sigma\,u(\lambda) \quad\text{i.e.}\quad \sigma(g)\,g^{-1}=u(\lambda)
\]
ou encore, \struck{en écrivant \ill{} de $g$} \struck{(puisque $g\in T_0$)}
\[
(141)\qquad \sigma(g)=u(\lambda)\,g
\quad \text{\struck{avec $\lambda=\delta(g)$}, pour $g\in T_0$ et $\lambda\in Z$ conv.}
\]
\note{au-dessus de $\sigma(g)$, un trait et un signe peu lisible ; au-dessus du second $g$, « $T_0\ni$ ».}
\struck{soit encore, posant $\tilde\delta(g)=\tilde\lambda$, $u_1(\tilde\lambda)=u(\lambda)\,u_0(\tilde\lambda)$}

\struck{Or a \ill{}} \add{En effet}, les deux membres sont dans $T_G$ (qui est normalisé par
$\sigma$), et \ill{} qu'ils \ill{} normalisant $U_0$, donc \ill{} ;
$T_G\cap U_0=T_H\cap U_0=\{1\}$, donc ils sont égaux. De plus, \struck{en posant cette
formule \ill{}} \add{cette formule montre que} $\sigma(g)$ \add{\uncertain{et $g$}}
\ill{} $= $ image de $\tilde\lambda$ dans $\widetilde{Z}=T_G/T_H$, donc,
\struck{\ill{}} ils \ill{} dans $T_0$ et $T_1$ respectivement, ils peuvent s'écrire
$u_0(\tilde\lambda)$, $u_1(\tilde\lambda)$, et (141) s'écrit
\[
(142)\qquad u_1(\tilde\lambda)=u(\lambda)\,u_0(\tilde\lambda)
\]
qui montre que
\[
(143)\qquad \lambda=c(\tilde\lambda) \quad (\text{cf }(119)).
\]

\page{264}
\note{p.~135 de l'auteur.}
En résumé :
\[
(144)\qquad
\left\{
\begin{array}{l}
T_0 \text{ opère fidèlement sur } H, \text{ et via cette opération s'identifie}\\
\text{au sous-groupe } \widetilde{Z} \text{ de } \mathrm{Aut}(H) \text{ formé des } \bar g\in\mathrm{Aut}\,H\\
\text{satisfaisant les conditions}\\
\qquad \begin{cases}
\bar g \text{ normalise } U_0 \text{ et centralise } T_H\\
\bar g(\sigma)\in\sigma T_H .
\end{cases}\\
\text{De plus, on aura alors}\\
\qquad \bar g(\sigma)=\sigma\,u(\delta(g)) \qquad \delta(g)=c(\tilde\delta(g)) .
\end{array}
\right.
\]
\note{la mention « $\widetilde{Z}$ » est ajoutée au-dessus de « cette » ; l'égalité $\delta(g)=c(\tilde\delta(g))$ est écrite sous $\delta(g)$, reliée par une accolade.}
\marginal{NB Ces conditions impliquent aussi que $\bar g$ normalise $U_1$, --- en fait, si $\bar g$ \struck{\ill{}} vérifie $\sigma T_H=T_H\sigma$ \ill{} \ill{}, alors il normalise $U_1$.}

En d'autres termes, \struck{si on ne distingue pas \ill{} $\widetilde{Z}$ \ill{}}
\add{« copie conforme » de $T_0\subset\mathrm{Aut}(H)$, via $u_0\colon\widetilde{Z}\xrightarrow{\sim}T_0$ et
$\tilde\delta_{T_0}\colon T_0\xrightarrow{\sim}\widetilde{Z}$,}
alors l'hom.\ $c\colon\widetilde{Z}\to Z$ s'identifie à l'hom.
\[
(145)\qquad \widetilde{Z}\xrightarrow{c_0}T_H
\]
défini par
\[
\bar g(\sigma)=\sigma\,c_0(\bar g)
\]
\note{le passage « copie conforme … $\widetilde{Z}$ » est encadré et renvoyé par des traits ; l'ordre retenu ici est une lecture. À droite, séparé par un trait vertical :}
(145$'$) i.e.\ on a commutativité dans

\begin{tikzcd}
\widetilde{Z} \arrow[r, "c"] \arrow[dr, "c_0"'] & Z \arrow[d, "u", "\wr"'] \\
 & T_H
\end{tikzcd}

\note{sous $\widetilde{Z}$, une flèche verticale $u_0$ vers $T_0$, et de là une flèche vers $T_H$, sont biffées en zigzag.}

Notons d'autre part que l'hom.
\[
T_H\longrightarrow\mathrm{Aut}(H),\qquad t\longmapsto \operatorname{int}_H(t)
\]
se factorise visiblement par $\widetilde{Z}$, et donne un homomorphisme
\[
(146)\qquad T_H\xrightarrow{c_0'}\widetilde{Z}
\]
satisfaisant
\[
(147)\qquad c_0(c_0'(t))=t^{-2}
\]
puisque pour $t\in T_H$, on a $\sigma t\sigma^{-1}=t^{-1}$ \struck{\ill{}} i.e.\
$t\sigma t^{-1}=\sigma t^{-2}$, ce qui suggère qu'on a commutativité dans

\page{265}
(148)

\begin{tikzcd}
\widetilde{Z} & Z \arrow[l, "c'"'] \arrow[d, "u", "\wr"'] \\
 & T_H \arrow[ul, "c_0'"]
\end{tikzcd}

\[
\operatorname{int}_H(t^{-1})=c_0'(t^{-1})\longleftarrow\!\!\dashv\ t
\]
\note{le libellé de la diagonale est surchargé ; une flèche verticale $u_0$ sous $\widetilde{Z}$ et un trait horizontal vers $T_H$ sont biffés en zigzag, comme p.~264.}

i.e.\ qu'on doit avoir
\[
\operatorname{int}_H(u(\lambda))^{-1}\ \text{\struck{\ill{}}}\ =\operatorname{int}(u_0(c'(\lambda)))|H
\]
ou encore
\[
u(\lambda)^{-1}\equiv u_0(c'(\lambda)) \mod \operatorname{Cent}_G(H)
\]
\struck{i.e.\ que}
\[
u(\lambda).u_0(c'(\lambda))\in\operatorname{Cent}_G(H)
\]
or on a
\[
(128)\qquad u(\lambda)\,u_0(c'(\lambda))=\lambda\in Z\subset\operatorname{Cent}(G),
\]
ce qui \struck{\ill{}} prouve (148).

Cela achève donc d'expliciter ce qu'il fallait via $T_0$. On aurait l'explicitation
symétrique via $T_1$, avec $U_1$ remplaçant $U_0$
\struck{--- en fait on trouve \ill{} le même sous-groupe \ill{} $T_G$, \ill{} façon \ill{} --- mais en
fait $T_G$ n'opère pas \emph{fidèlement} sur $H$, donc $T_G$ \ill{} s'explicite comme un
sous-groupe de $\mathrm{Aut}\,H$.}
\[
(149)\qquad \widetilde{Z}=\operatorname{Im}(T_0\to\mathrm{Aut}(H))=\operatorname{Im}(T_1\to\mathrm{Aut}(H))
=\operatorname{Im}(T_G\to\mathrm{Aut}\,H)
\]
\note{ce passage, rayé par de grands traits obliques et encadré, porte l'égalité (149) écrite en surcharge ; seule celle-ci semble maintenue.}
\struck{Mais} Explicitons encore le sous-groupe $\mathfrak{z}$ \add{$=\operatorname{Ker}(\delta_{T_0}\colon T_0\to Z)$}
de $T_0$ en tant que sous-groupe de $\mathrm{Aut}\,H$, on trouve
\[
(149)\qquad \mathfrak{z}\xrightarrow{\ \sim\ }
\left\{\begin{array}{l}
\text{sous-groupe de } \mathrm{Aut}\,H \text{ formé des}\\
\text{automorphismes } \bar g \text{ qui normalisent } U_0,\\
\text{et qui centralisent } T_H \text{ et } \sigma.
\end{array}\right.
\qquad z\mapsto \operatorname{int}(z)|H
\]
\struck{et aussi $\sigma(U_0)$} qui
\note{le numéro (149) est donné deux fois ; le $T_H$ de la dernière ligne est surchargé et se lit mal ; la lettre $\mathfrak{z}$ est un z ronde.}

\page{266}
\note{p.~136 de l'auteur.}
\emph{NB} Introduisons
\[
(150)\qquad N_H=T_H\cup\underbrace{(T_H\sigma=\sigma T_H)}_{\Pi=T_H\sigma=\sigma T_H}
\]
qui est un sous-groupe de $H$ \struck{normalisant} contenant $T_H$ et le normalisant, avec $T_H$
d'indice 2 dans $N$, ou égal à $N$ dans le cas $\sigma\in T_H$, \struck{d'ici} cas \ill{}
que nous allons écarter pour le moment (cas « \uncertain{non} dégénéré ») :
\[
(151)\qquad 1\longrightarrow T_H\longrightarrow N_H\longrightarrow\{\pm1\}\longrightarrow 1
\qquad \text{si } \sigma\notin T_H
\]
On a donc
\[
(152)\qquad \widetilde{Z}=\left\{\bar g\in\mathrm{Aut}(H)\;\middle|\;
\begin{array}{l}
g \text{ centralise } T_H \text{ et normalise } N_H\\
g \text{ normalise } U_0 \text{ (ou encore : } U_1)
\end{array}\right\}
\]
\note{dans (152), « $T_H$ » est surchargé en « $N_H$ » ou l'inverse, et « normalise » est écrit au-dessus d'un « centralise » biffé.}
Donc $\widetilde{Z}$ opère sur la classe $\Pi=\sigma T_H=T_H\sigma$ par translations, via
l'hom.\ $c\colon\widetilde{Z}\to Z$ justement.

\note{le paragraphe suivant, encadré, est rayé de grands traits obliques.}
\struck{Le cas « dégénéré » s'explicite ainsi. Conditions équivalentes :}
\[
(153)\qquad
\begin{cases}
\text{a) } t_0=s_0\\
\text{b) } A_0 \text{ est formé des parties ponctuelles de } S\\
\text{c) } \sigma=1\\
\text{d) } \sigma\in T_H \text{ i.e. } T_H=N_H\\
\text{e) } \ldots
\end{cases}
\]
\note{dans (153), c) et e) sont biffés ; e) commençait par « $S$ est un torseur sous $H$ ».}
\struck{Sous ces conditions on a $U_0=U_1=\{1\}$, $S$ est un torseur sous $H$, et $T_H$ est un
sous-groupe commutatif tel que $t\in T_H\Rightarrow t^2=1$, \ill{} enfin}
\struck{$\widetilde{Z}=\operatorname{Cent}_H(Z)$.}

Il y a lieu d'introduire aussi
\[
(154)\qquad N_G=T_G\cup(\sigma T_G=T_G\sigma)
\]
et on a donc, dans le cas où $\sigma\neq1$ (cas non dégénéré)
\[
(155)\qquad 1\longrightarrow T_G\longrightarrow N_G\longrightarrow\{\pm1\}\longrightarrow 1
\qquad \text{si } \sigma\notin T_H
\]
\note{les numéros (154) et (155) sont surchargés et se lisent mal ; le signe devant « 1 » dans « $\sigma\neq1$ » l'est aussi.}

\page{267}
Pour le cas « dégénéré », on trouve
\[
(155)\qquad
\begin{array}{c}
\sigma\in T_H \Longrightarrow U_0=U_1=\{1\},\ T_H \text{ un groupe commutatif avec } t\in T_H\Rightarrow t^2=1,\\
\text{enfin } \widetilde{Z}=\operatorname{Cent}_H(T_H).\\
\Updownarrow\\
\text{\struck{\ill{}, $t_0=\lambda s_0$}}\\
t_0\in Z.s_0\\
\text{i.e. } \underbrace{t_0^{\Diamond}}_{D_0}=\underbrace{s_0^{\Diamond}}_{D_1}
\end{array}
\]
\note{le signe en exposant, ici rendu $^{\Diamond}$, est un petit losange ou une boucle que nous ne savons pas identifier ; il revient plus bas ($S^{\Diamond}$).}
\marginal{implique donc que $S$ est un torseur sous $H$ ($\Leftrightarrow U_0=\{1\}$, $U_1=\{1\}$),}
\struck{Notons par là}
Plus précisément, les données d'un multigraphe : \uncertain{homothétie} et groupe simplement
transitif \uncertain{dans} \struck{\ill{}} \ill{} par $(s_0,t_0)\in A_0$, et tel que
\struck{tel que $\sigma\in T_H$ i.e.} $s_0^{\Diamond}=t_0^{\Diamond}$ i.e.\ $\sigma\in T_H$,
\emph{équivaut} : à la donnée de $H$, d'un sous-groupe commutatif $T_H$ tel que $T_H^2=\{1\}$, et
de $\sigma\in T_H$ --- et on a alors
\[
\widetilde{Z}=\operatorname{Cent}_H(T_H).
\]
\struck{Conditions} \add{Notons} le cas, encore plus dégénéré
\[
(156)\qquad \sigma=1 \iff s_0=t_0 \iff A_0 \text{ est formé des parties ponctuelles de } S .
\]


Digression dans le cas d'un multigraphe \add{de $Z$-}\uncertain{orientable} (de $Z$-\uncertain{homothétie}), et groupe
d'automorphismes $H$ simplement transitif sur \struck{les arêtes} \add{un des} $\vec{A}_\lambda$
(\uncertain{les autres} \ill{}). Fixons nous une origine $s_0$ et une arête incidente :
\[
(157)\qquad s_0\in S,\ D_1\in S^{\Diamond},\quad D_0=Z.s_0 \text{ et } D_1 \text{ incidents dans } S^{\Diamond},
\]
\note{à gauche de (157), « \struck{donc $s_0$,} » biffé ; « ch. une origine » au-dessus.}
donc
\[
(158)\qquad S\simeq H/U_0 \qquad U=H_{s_0} .
\]
\uncertain{ci-dessous} de $H$-\uncertain{ensembles}),

\struck{Fixons nous de plus un \ill{}} De plus $T_G$, $T_0$ \struck{\ill{}} $\subset\mathfrak{S}_S$ sont bien
déterminés par (137), \struck{et $N_1$} \ill{} (90) (97), donc aussi
$T_H=T_G\cap H$ :
\struck{$T_H\subset H$, $T_0\subset T_G\subset G$}

\page{268}
\note{p.~137 de l'auteur.}
(159)

\begin{tikzcd}[column sep=small, row sep=small]
 & e \arrow[r, hook] \arrow[d, hook] & U_0 \arrow[d, hook] \\
 & T_H \arrow[r, hook] \arrow[d, hook] & H \arrow[d, hook] \\
T_0 \arrow[r, hook] & T_G \arrow[r, hook] & G
\end{tikzcd}

\note{les inclusions sont écrites $\subset$ sur la page, verticales et horizontales.}

Bien sûr, $\sigma$ n'est pas déterminé en termes des données précédentes, mais on a
\[
(160)\qquad
\begin{cases}
N_G=\{g\in G\mid u(\{D_0,D_1\})=\{D_0,D_1\}\},\\
N_H=N_G\cap H=\{g\in H\mid u(\{D_0,D_1\})=\{D_0,D_1\}\}
\end{cases}
\]
\note{« \struck{donc $N_H$ \ill{}} » biffé devant la seconde ligne ; le $u$ des accolades est sans doute pour $g$.}
d'où l'hom.
\[
N_G\longrightarrow\{\pm1\} \qquad (\text{si } \sigma\notin H \text{ i.e. } D_0\neq D_1)
\]
étant donné par l'action de $g$ sur la paire $D_0,D_1\ldots$

Ainsi $N_G$, $N_H$ et l'hom.\ $N_G\to\pm1$ \struck{\ill{}} du noyau $T_G$ (resp.\ $N_H\to\pm1$ du
noyau $T_H$) sont déterminés par les données. Ceci dit, on a des \struck{bijections} \add{isom.}\
canoniques de $Z$-torseurs
\[
(161)\qquad \text{\struck{\ill{}}}\quad
\underbrace{N_H^{-}}_{\substack{\text{image inverse de } -1\\ \text{par } N_H\to\{\pm1\}}}
\xrightarrow{\ \sim\ }D_1\xrightarrow{\ \sim\ }\Pi
\qquad \sigma\mapsto\sigma(s_0),\quad t_0\mapsto\tau(s_0,t_0)
\]
\note{devant $N_H^{-}$, un premier essai est biffé en zigzag ; sous l'accolade, un « $N_H\setminus H$ » noirci.}
\struck{\ill{}} où on prend sur $N_H^{-}$ la structure de $Z$-torseur définie par les opérations
\emph{droites} de $T_H\xleftarrow[u]{\sim}Z$ sur $N_H^{-}$, i.e.
\[
(162)\qquad \lambda\in Z \text{ opère sur } N_H^{-} \text{ par } g\longmapsto g\,u(\lambda) .
\]
Notons ici :

\page{269}
\[
(163)
\]
Il y a une \struck{\ill{}} \add{équivalence} entre les données $(S,H,Z,\Pi,s_0,D_1)$ d'un
multigraphe \add{simple} : homothétie de $Z$-\uncertain{orbites}, avec groupe d'automorphismes $H$
transitif sur les $\vec{A}_\lambda$, un sommet marqué $s_0$ et une \uncertain{voie} $D_1$ incidente :
$Zs_0=D_0$ \uncertain{marquée}, \add{$D_0\neq D_1$ ($\Leftrightarrow$ cas « non dégénéré »)}
\struck{d'une part},

et d'autre part : les systèmes $(H,U_0,T_H,N_H)$ d'un groupe $H$, et de ss-groupes $U_0$, $T_H$,
$N_H$, satisfaisant :
\begin{itemize}
\item[a)] $T_H$ normalise $U_0$, et $T_H\cap U_0=\{1\}$ (d'où produit semi-direct $B_H=T_H.U_0$)
\item[b)] $N_H\supset T_H$, $T_H$ d'indice 2 dans $N_H$ (qui normalise donc $T_H$) --- de sorte
que les $\sigma(U_0)$ pour $\sigma\in N_H\setminus T_H=N_H^{-}$ \struck{\ill{}} sont égaux : un
même ss-groupe $U_1$ de $H$
\item[c)] \struck{$T_H$ commutatif et $N_H/T_H\simeq\{\pm1\}$ opère sur $T_H$ par}
$\sigma\in N_H^{-}\Longrightarrow\sigma^2=1$ \struck{\ill{}}
($\Rightarrow T_H$ est commutatif et $N_H/H\simeq\pm1$ y opère par $t\mapsto t^{-1}$)
\item[d)] $U_0\cap U_1=\{1\}$ (où $U_1$ défini dans b)).
\end{itemize}
\note{dans c), « $N_H/H$ » est sans doute pour $N_H/T_H$.}

On aimerait vérifier aussi la formule familière
\[
B_H\cap N_H=T_H
\]
\marginal{donc les données $(U_0,T_H,N_H)$ équivalent \ill{} aux $(B_H,N_H)$, \ill{} (164). \ill{} $T_H=B_H\cap N_H$ ; \ill{} $U_0=[B_H,B_H]$ \ill{} et \ill{} déterminés par $B_H$, $N_H$ de $H$.}
\note{note écrite en oblique dans la marge de gauche ; lecture très incertaine.}
i.e.\ le groupe $B_H\cap N_H$ est connu \uncertain{entre} $T_H$ et $N_H$ --- donc si
$\sigma\in N_H^{-}$ est donné, la formule équivaut à
\[
(164)\qquad \sigma\notin B_H
\]
\[
(165)\qquad
\]
Mais si \uncertain{on avait} $\sigma\in B_H$
\struck{$=B_H\,U_0\,T_H$, $B_H=T_H$ donc $\sigma\in T_H$, donc $u_0\in U_0$, $t\in T_H$, donc
\ill{} centralise $T_H$}
$\sigma=ut$ \ldots
\note{le reste de la page, encadré, est rayé de quatre grands traits obliques ; on en donne la lecture :}
\struck{$\sigma(U_0)=U_1$ \ill{}. Il normalise et y induit $g\mapsto g^{-1}$, donc il faudrait que $u\in U_0$ normalise $T_H$ et y induit $g\mapsto g^{-1}$ ; cela n'est possible, comme $T_H$ normalise $U_0$ et $T_H\cap U_0=\{1\}$, que si $g\in T_H\Rightarrow g^2=1$ i.e.\ $g=g^{-1}$. Cela veut dire que $B_H\cap N_H=N_H$ i.e.\ $N_H\subset B_H$, et alors $T_H^2=\{1\}$ et $u$ centralise $T_H$). Donc (164), \ill{} bien}
\[
\text{\struck{$(165)\quad N_H=\{1,u_0\}\times T_H$}}
\]
\struck{où $u_0\in U_0$ est un élément d'ordre 2 qui centralise $T_H$, i.e.\ $u_0\in{}_2(U_0^{T_H})$, $u_0\neq1$.}

\note{la page 269 s'arrête sur cette discussion ; la p.~270 ouvre un nouveau paragraphe.}

\page{270}
\note{p.~138 de l'auteur.}
\subsection*{4. Cas des graphes triangulants}

Un graphe $(S,A_0)$ \add{sans boucles \struck{isolées}} est dit \emph{triangulant} si
\[
(1)\qquad \forall s,t\in S,\ \text{\struck{tels que}}\ \{s,t\}\in A_0,\ \exists u\in S \text{ tel que } \{s,u\},\{t,u\}\in A_0
\]
i.e.\ tel que $\{s,t,u\}\in F_0$, où on a posé
\[
(2)\qquad F_0=\{f\in\mathfrak{P}_3(S)\mid \forall a\in\mathfrak{P}_2(f),\ \text{on a } a\in A_0\} .
\]
\struck{La condition d'homogénéité \ill{}}
Lorsqu'on est en présence d'un \emph{multigraphe} : homothétique de $Z$, on dit qu'il est
triangulant si chacun des structures de graphes \add{\uncertain{induites}} est triangulant. Il suffit,
lorsque les \struck{\ill{}} structures sont isomorphes entre elles (« condition d'isomorphie » du
§~2) que l'une d'elles soit triangulante. Mais il n'est pas dit pour \uncertain{nous} qu'il en soit
ainsi dans tous les cas.

\page{272}
\note{fiche de petit format, écrite en travers ; un tableau.}

\[
\begin{array}{l||c||c|c|c}
n & 2 & 3 & 4 & 5\\ \hline
[\mathrm{SL}(2,\mathbb{Z}/n\mathbb{Z})] & 6 & 24 & 48 & 120\\
{}[\mathrm{GL}(2,\mathbb{Z}/n\mathbb{Z})] & 6 & 48 & 96 & 480\\ \hline
{}[\mathrm{SL}(2,\mathbb{Z}/n\mathbb{Z})'] & 6 & 12 & 24 & 60\\
\mathrm{SL}(2,\mathbb{Z}/n\mathbb{Z})' & \mathfrak{S}_3=D_3 & \mathfrak{S}_4^{+} & \mathfrak{S}_4 & \mathfrak{S}_5^{+}\\ \hline
{}[\mathrm{GP}(1,\mathbb{Z}/n\mathbb{Z})] & 6 & 24 & 48 & 120\\
\mathrm{GP}(1,\mathbb{Z}/n\mathbb{Z}) & \mathfrak{S}_3=D_3 & \mathfrak{S}_4 & \mathfrak{S}_4\times\{\pm1\}\,? & \mathfrak{S}_5\,?\\ \hline
{}[M^{*}] & 3 & 8 & 12 & 24\\
{}[\underbrace{M^{*}/\{\pm1\}}_{S_M}] & 3 & 4 & 6 & 12\\ \hline
\text{type } M^{*}/\{\pm1\} & \text{trièdre} & \text{tétraèdre} & \text{octaèdre} & \text{icosaèdre}\\ \hline
{}[P(M)] & 3 & 4 & 6 & 6\\ \hline
{}[F_M^{*}] & 1 & 4 & 8 & 40\\
{}[F_M^{\circ}] & 1 & 4 & 8 & 20
\end{array}
\]
\note{la notation $\mathrm{GP}(1,\ldots)$ est une lecture (« GP » ou « GPL ») ; dans la case « trièdre », un premier mot est noirci ; dans la ligne $[F_M^{*}]$, les chiffres 4, 8 et 40 sont surchargés sur d'autres, et l'en-tête de cette ligne porte un ajout noirci ; la ligne $[F_M^{\circ}]$, plus petite, est écrite dessous en correction.}

Soit $M$ un module libre de rang 2 sur \ill{}
\note{la fiche est déchirée en bas ; on y devine encore « \ldots dérons ».}

\page{273}
\note{verso de la fiche précédente : seule inscription, en haut à gauche.}
\[
(\mathbb{Z}/n\mathbb{Z})^{2*}
\]

\page{275}
\subsection*{Détermination d'un sommet (d'un \add{2-}polyèdre triangulé) par son étoile}

1) Soit \add{$P=$} $(S,F\subset\mathfrak{P}_3(S))$ un 2-polyèdre triangulé convexe, $s,s'\in S$
tels que \struck{\ill{}} l'on ait
\[
s\neq s',\qquad \mathrm{Et}(s)=\mathrm{Et}(s') .
\]
Alors on a
\[
S=\{s,s'\}\cup\mathrm{Et}(s)
\]
i.e.\ $P$ s'identifie au double cône sur le $n$-polygône $\mathrm{Et}(s)$
($n\geqslant3$) \add{$n=\operatorname{card}\mathrm{Et}(s)$}.

\drawing{à gauche du NB, un double cône (bipyramide) sur un pentagone, arêtes cachées en pointillé.}

\emph{NB} Les deux sommets $s,s'$ sont d'ordre $n$, les autres sommets sont d'ordre 4. Si tous les
sommets ont même ordre, on a donc $n=4$, et $P$ est un octaèdre. Dans ce cas, pour que deux
sommets aient même étoile, il faut et il suffit qu'ils soient antipodiques. Dans le cas $n\neq4$,
la seule paire de sommets ayant même étoile est $\{s,s'\}$ (car \struck{\ill{}} s'il y avait
une \struck{\ill{}} autre telle paire $\{t,t'\}$, \struck{les} l'un de $t,t'$ serait
\uncertain{quatre}, et $P$ serait un octaèdre, absurde !). Mais dans ce cas, il existe encore un
automorphisme unique $\underline{a}\neq\mathrm{id}$, (qui mérite le nom d'antipodisme ?) \uncertain{i.e.} tel
que pour \struck{\ill{}} $t\in S$, on ait $\mathrm{Et}(t)=\mathrm{Et}(a(t))$ ;
cet automorphisme est la transposition de $s,s'$ et l'identité sur
$E=\mathrm{Et}(s)=\mathrm{Et}(s')$. Il est involutif, et ses orbites sont les classes pour la
relation

\page{276}
d'équivalence \add{dans} \struck{\ill{}} $S$
\[
t\sim t' \overset{\text{déf}}{\iff} \mathrm{Et}(t)=\mathrm{Et}(t') .
\]
Cette dernière propriété est une caractérisation de $a$ commune à tous les cas (aussi bien $n=4$
que $n\neq4$).

2) Soit $\Lambda$ un anneau semi-local, \struck{\ill{}} \add{$M$ un} module libre de rang 2 sur
$\Lambda$, $u,u'\in M^{*}$ tels que l'on ait
\[
(*)\qquad
\forall x\in M^{*},\quad u\wedge x\in(\det M)^{*}
\ \overset{\text{\struck{$\iff$}}}{\Longrightarrow}\
\text{\struck{$u'\wedge x\in(\det M)^{*}$}}\ \text{\add{et}}\ u'\wedge x\in\pm\{u\wedge x\}
\]
Alors, on a $u'\in\{\pm u\}$, sauf dans le cas suivant :
\struck{$2\in\ldots$} $2\Lambda=\{0,2\}\neq\{0\}$ \struck{\ill{}} i.e.
\[
\left\{
\begin{array}{l}
\Lambda \text{ une } \mathbb{Z}/4\mathbb{Z}\text{-alg.}\ \text{\struck{\ill{}}}\ \text{est extension}\\
\text{d'une } \mathbb{F}_2\text{-algèbre } \Lambda_0=\Lambda/2\Lambda\\
\text{par un idéal de carré nul}\\
\text{engendré par } 2_\Lambda \text{ et qui est}\\
\text{de rang 1 en tant que } \mathbb{F}_2\text{-}\\
\text{module [correspondant donc à}\\
\text{une \emph{connexion} } \chi\colon\Lambda_0\to\mathbb{F}_2]
\end{array}
\right.
\]
\marginal{NB le cas exceptionnel est aussi celui où le $\mathrm{GL}_1(M)^{\pm}$ est normalisant dans $\mathfrak{S}_{S_M}=\mathfrak{S}(M^{*}/\pm1)$, \ill{} \ill{} \ill{} \ill{} les groupes induisent opérations \ill{} $\mathrm{Aut}/(\pm1)$ \ill{} \ill{} \ill{} podismes \ill{}.}
\note{note écrite en oblique dans la marge de gauche ; lecture très incertaine.}

Notons d'abord \struck{\ill{}} \emph{lemme} : La \struck{condition} \add{relation}
« $\forall x\in M^{*}$, $u\wedge x\in\det M^{*}\Longrightarrow u'\wedge x\in\det M^{*}$ »
équivaut à « $\forall\mathfrak{m}_i\in\operatorname{Max}(\Lambda)$, \struck{\ill{}} les images
$u_i^{\circ}$ et $u_i'^{\circ}$ de $u,u'$ dans le corps résiduel $k_i=\Lambda/\mathfrak{m}_i$ sont
\struck{proportionnelles} \add{colinéaires} »
[\emph{NB} les conditions $u,u'\in M^{*}$ \struck{\ill{}} impliquent que $u_i^{\circ}$, $u_i'^{\circ}$
sont $\neq0$] --- c'est \add{donc} une relation d'\emph{équivalence} entre $u,u'\in M^{*}$. Cette
relation implique l'existence d'une base $e,f$ de $M^{*}$, telles que l'on ait
$u\wedge e,\ u\wedge f,\ u'\wedge e,\ u'\wedge f\in\det M^{*}$, i.e.
\[
u=ae+bf,\quad u'=a'e+b'f\qquad \text{avec } a,b,a',b'\in\Lambda^{*}
\]
et, \uncertain{quitte} à remplacer $e,f$ par $ae$, $bf$, qu'on a \uncertain{simplement}
$a=b=1$, i.e.\ $u=e+f$.
Si on suppose maintenant
$u\wedge x\in\det M^{*}\Longrightarrow u'\wedge x\in\pm\{u\wedge x\}$, faisant $x=e$, puis $x=f$,
on voit que \struck{$a'\in\{\pm a\}$}, $a'\in\{\pm1\}$, $b'\in\{\pm1\}$, donc quitte à changer
$u$ de signe, $u=e+f$, $u'=e\pm f$, ops

\page{277}
et il s'agit de voir que l'on ne peut avoir
\[
u=e+f,\quad u'=e-f
\]
sauf \struck{dans le cas précisé} si $2\cdot1_\Lambda=0$ (\uncertain{alors} $+1_\Lambda=-1_\Lambda$ donc
$u=u'$, et \struck{on} dans le cas précisé dans l'énoncé ci-dessus.) \struck{Supposons donc}
\[
2\cdot1_\Lambda\neq0 .
\]
Notons que \struck{puisque $u\equiv u'$} \add{comme $u_i^{\circ}$, $u_i'^{\circ}$ sont colinéaires}
dans le corps résiduel $k_i$, cela implique déjà que $2\cdot1_{k_i}=0$ pour tout $i$, i.e.
\[
2\cdot1_\Lambda\in\mathfrak{r}\Lambda .
\]
\struck{La} condition $(*)$, sur $x\in M^{*}$ s'écrit, écrivant
\[
x=\alpha e+\beta f\qquad \alpha,\beta\in\Lambda
\]
que
\[
\alpha+\beta\in\Lambda^{*}\Longrightarrow \text{\struck{\ill{}}}\ \alpha-\beta\in\pm(\alpha-\beta)
\quad\text{i.e.}\quad 2\alpha \text{ ou } 2\beta \text{ nul.}
\]
\note{la conclusion porte sans doute sur $\alpha+\beta\in\pm(\alpha-\beta)$ ; nous transcrivons ce que porte la page.}
Posant $\lambda=(\alpha+\beta)$ et multipliant $\alpha,\beta$ par $\lambda^{-1}$, on est ramené à voir que
\[
\alpha,\beta\in\Lambda,\ \alpha+\beta=1\Longrightarrow 2\alpha \text{ ou } 2\beta \text{ nul \struck{\ill{}}, i.e.}
\]
$\forall\alpha\in\Lambda$ on a $2\alpha$ ou $2(1-\alpha)=2-2\alpha$ nul, i.e.
\[
2\Lambda\subset\{0,2\},\ \text{donc (puisque } 2\cdot1_\Lambda\neq0)
\]
\note{à gauche, « (\struck{\ill{}}) » biffé.}
\[
\boxed{2\Lambda=\{0,2\}\neq\{0\}}\ ,\quad (\text{qui implique } 4_\Lambda=0\ \ldots)
\]
\emph{Remarque} (On voit que dans ce cas, les \struck{couples} \add{paires} $\{u,u'\}$ qui mettent
en défaut la conclusion $\{u'\in\pm u\}$ forment une orbite sous $\mathrm{GL}_\Lambda(M)$, \struck{\ill{}}
on peut aussi écrire, plus précisément
\[
u'=u-2f=u+2f\quad(\text{NB } 4_\Lambda=0),
\]
\[
-u'=-u-2f=-u+2f)
\]
où $2f$ ne dépend que de $f$ modulo l'idéal $\mathfrak{m}_{i_0}$ annulateur de $2_\Lambda$, qui est un
idéal maximal : corps résiduel $k_{i_0}\simeq\mathbb{F}_2$ \struck{$\mathbb{F}_2$}
\[
\text{on a \struck{\ill{}} une suite exacte}\quad
0\longrightarrow\underbrace{M\otimes_\Lambda k_{i_0}}_{\substack{\simeq 2M\\ \text{vectoriel de}\\ \text{rang 2 sur } k_{i_0}\simeq\mathbb{F}_2}}\longrightarrow M\longrightarrow M/2M\longrightarrow 0
\]

\page{278}
Ainsi on peut écrire
\[
u'=u+\xi=u-\xi\qquad(\text{car } 2\xi=0)
\]
où $\xi$ est un des trois éléments $\neq0$ de $2M\simeq M\otimes_\Lambda k_{i_0}$, et plus
précisément, un des deux éléments \struck{\ill{}}
\[
\xi\in\bigl[(2M=M\otimes_\Lambda k_{i_0})\setminus\{0,2u\}\bigr]
\]
(car $\xi\notin\{0,2u\}$ signifie justement que $u+\xi\notin\{\pm u\}$).
Si
\[
2M\setminus\{0,\underset{?}{2u}\}=\{\xi',\xi''\},
\]
notons que alors
\[
\underbrace{\text{\struck{\ill{}}}\ \eta}_{2u}+\xi'+\xi''=0
\]
donc
\[
u+\xi'=-(u+\xi'')\quad \text{\struck{\ill{}}}\qquad -(u+\xi'')=-u-\xi''
\]
\note{la seconde expression, $-u-\xi''$, est écrite sous une accolade placée sous $-(u+\xi'')$.}
i.e.\ l'image de $u'$ dans $M^{*}/\{\pm\mathrm{id}_M\}=S_M$ est uniquement déterminée par $u$.
Ainsi on peut ainsi \uncertain{définir} une bijection canonique
\[
\underline{a}\colon S_M\longrightarrow S_M
\]
qui associe : à tout $\dot u\in S_M$ l'unique $\dot u'\in S_M$ tel que l'on ait
\[
\mathrm{Et}_i^{*}(\dot u)=\mathrm{Et}_i(\dot u'),\qquad \dot u'\neq\dot u
\]
pour un $i\in\pi_M=(\det M)^{*}/\pm1_\Lambda$ --- \uncertain{et} qui est alors vrai pour tous.
On voit de suite que $\underline{a}$ est involutif, \add{et} sans pts fixes par construction.
On peut \struck{\ill{}} \add{décrire} \struck{\ill{}} antipodismes par la suite exacte
\[
0\longrightarrow \underbrace{\text{\struck{\ill{}}}\ N}_{\substack{\simeq\,\mathrm{End}_{\mathbb{F}_2}(M_0)\\ \text{où } M_0=M\otimes_\Lambda(k_{i_0}\simeq\mathbb{F}_2)}}
\hookrightarrow \mathrm{Aut}_\Lambda(M)\longrightarrow\mathrm{Aut}_{\Lambda/2\Lambda}(M/2M)\longrightarrow1
\]
et on a
\[
\mathrm{End}(M_0)\simeq M_0\oplus\mathbb{F}_2^{\,\underline{\omega}_J}\qquad(J=M_0^{*}) ,
\]
Donc
\[
\text{\struck{\ill{}}}\ N/\{\pm\mathrm{id}_M\}\simeq M_0\oplus\mathbb{F}_2^{\,\underline{\omega}_J}/\mathrm{diag}\simeq M_0+\mathbb{F}_2
\]
et l'élément non nul du facteur $\mathbb{F}_2$ donne $\underline{a}$.
\note{la lettre notée ici $N$ est une lecture ; « $\mathbb{F}_2^{\,\underline{\omega}_J}$ » aussi : un exposant souligné, peut-être $\underline{\omega}_J$ ou $\simeq J$.}

\page{279}
\struck{Le cas} Le cas le plus simple est celui où
\[
\Lambda=\mathbb{Z}/4\mathbb{Z}
\]
c'est aussi celui où $(\text{\struck{\ill{}}}\ S_M,F_M)$ est justement l'octaèdre, dont on retrouve
l'antipodisme \add{particulier} \struck{\ill{}} d'algèbre linéaire. La construction de
l'antipodisme est évidemment indépendante de l'hyp.\ semi-locale sur $\Lambda$, elle marche dès que
\[
2\Lambda=\{0,2\}\neq\{0\},
\]
et dans $S_M$ on a encore
\[
\begin{cases}
\mathrm{Et}_i(u)=\mathrm{Et}_i(\underline{a}(u)) & \forall u\in S_M,\ i\in\pi_M\\
\underline{a}(u)\neq u
\end{cases}
\]
--- mais il n'est pas \uncertain{vrai} nécessairement qu'il n'y a pas d'autres éléments que
$u,u'=\underline{a}(u)$ \uncertain{ayant même} étoile que $u$ \ldots\ \struck{(cas dé\ldots)}

\emph{NB} Dans tous les cas, on a
\[
\underline{a}\in\mathrm{GL}'_\Lambda(M)^{\pm}
\]
car comme
\[
\underline{a}\equiv\mathrm{id}\mod 2\Lambda=\{0,2\},\ \text{on a}
\]
\[
\det\underline{a}\in1+\{0,2\}=\{\pm1\} .
\]
En fait, on vérifie que
\[
\det\underline{a}=-1
\]
(car les éléments de $N$ qui sont dans $\mathrm{SL}_\Lambda(M)$ sont ceux de
$\operatorname{Im}(M_0\oplus\mathbb{F}_2\,\mathrm{id})$, et $\underline{a}\in N\simeq\mathrm{End}_{\mathbb{F}_2}(M_0)$
n'est pas dans ce sous-module\ldots

\page{280}
\note{page écrite au crayon, en brouillon rapide.}
\[
\left\{
\begin{array}{ll}
1) & \{s,t\}\in A_0,\ \text{\struck{\ill{}}}\ \Longrightarrow\ \text{\struck{$\exists$}} \text{ existe exactement deux}\\
 & \text{éléments \struck{\ill{}} tels que } \{s,u\} \text{ et } \{t,u\}\in A_0\\
 & \text{\struck{$\{s,t\}$}}\\
2) & \forall s,t\in S,\ \exists u\in S \text{ tel que } \{s,u\},\{t,u\}\in A\\
 & \text{et OPS } \{s,u\}\in A_0
\end{array}
\right.
\]
\struck{Connexe} \emph{Condition} $\forall s,t\in S$, $\exists u,v,w\in S$ tels que
$(s,u,v,$ \struck{$w$}$,t)$ soit une $A_0$ chaîne.

\drawing{à gauche, un petit schéma : $s$ et $t$ reliés par des segments à $u$ et $v$, avec des petits cercles marquant des sommets ; dessous, un essai griffonné et biffé.}

\emph{Dém.} Par 2) \add{pour $(s,t)$} on déduit $u$ avec
\[
\{s,u\}\in A_0,\qquad \{u,t\}\in A
\]
puis par 1) pour $(u,t)$, on déduit $v$ qui relie $u$ à $t$ pour $A_0$

\drawing{un losange : $u$ à gauche, $v$ à droite, $s$ en haut, $t$ en bas ; la diagonale $uv$ et les côtés marqués $A$ et $A_0$ ; suivi de trois traits verticaux.}

3) \emph{$A(s)=A(t)\Longrightarrow\exists(u,v)\in S^2$, \ldots\ ?}

\emph{Sommet déterminé par son étoile ?}

Pas vrai pour $\Lambda=\mathbb{Z}/4\mathbb{Z}$ (octaèdre)

\drawing{dans la marge de gauche, au crayon léger, un octaèdre en perspective.}

\struck{\ill{}} en deux éléments de $M^{*}$, savoir $u$ et \struck{\ill{}} $\underline{a}(u)$
\struck{$=u+\ldots$, $u+2u=3u=\ldots=u+2v$ (\ldots)}
\[
\text{\struck{$u'=u$}}\quad \text{\struck{\ill{}}}\quad \text{\struck{$M^{*}_2=$}}\qquad x\wedge u=\pm x'\wedge u\quad\text{quand}
\]
\[
x\wedge u \text{ et } x'\wedge u \text{ sont }\in\Lambda^{*}
\]
\[
\begin{array}{lll}
x,\text{\struck{$x'$}}=\lambda e+\mu f & x\wedge e=\pm\mu & x\wedge f=\pm\lambda\\
x'=\lambda'e+\mu'f & x'\wedge e=\pm\mu' & x'\wedge f=\pm\lambda'\\
 & \text{OPS} & x=e+\text{\struck{\ill{}}}\,f \qquad u=\ldots\\
 & & x'=e-f
\end{array}
\]
\[
\text{\struck{$x\wedge$}}\ u=ae+bf
\]
\[
\begin{array}{l}
x\wedge u=-a+b\\
x'\wedge u=+a+b
\end{array}
\qquad \text{\struck{$a+b$}}\quad 2a=0\ \text{ou}\ 2b=0\quad \text{\struck{si $(a,b)\in\Lambda$}},\ b-a,\ b+a\in\Lambda^{*}
\]
\note{la page s'arrête sur ce calcul ; nous ne savons pas s'il se poursuit au-delà du lot.}

\end{document}
